\subsection{分式模}

第 1 节中定义的典范同态 \(i_{\mathbf{A}}^{\mathbf{S}} \colon \mathbf{A} \to \mathbf{A}[\mathbf{S}^{-1}]\) 使我们可以把每个 \(\mathbf{A}[\mathbf{S}^{-1}]\) -模看作 A-模。

命题 3. 设 A 是环，S 是 A 的子集，M 是 A-模，M' 是 A-模 M ⊗\_A A{[}S\^{}\{-1\}{]}，f 是 M 到 M' 的典范 A-同态 x ↦ x ⊗ 1。那么：

\begin{enumerate}
\def\labelenumi{(\arabic{enumi})}
\item
  对所有 \(s \in \mathbf{S}\)，位似映射 \(z \mapsto sz\) 在 \(\mathbf{M}'\) 上是双射。
\item
  对每个满足如下条件的 A-模 N：对所有 \(s \in S\)，位似映射 \(y \mapsto sy\) 在 N 上是双射；以及对每个 M 到 N 的同态 u，存在唯一的同态 \(u'\)（\(M'\) 到 N），使得 \(u = u' \circ f\)。
\end{enumerate}

换言之，\((M',f)\) 是泛映射问题（《集合论》第 IV 章，§ 3，第 1 节）在下列条件下的解：结构种 \(\Sigma\) 是这样的 A-模的结构种，其中由 S 的元素诱导的位似映射都是双射；态射是 A-模同态，而 \(\alpha\) -映射也是 A-模同态。

对每个 A-模 N 与所有 \(a \in \mathbf{A}\)，用 \(h_a\) 表示位似映射 \(y \mapsto ay\)（在 \(\mathbf{N}\) 上）；于是 \(a \mapsto h_a\) 是 A 到 \(\operatorname{End}_{\mathbf{A}}(\mathbf{N})\) 的环同态。说 \(h_a\) 是双射，就是说 \(h_a\) 是 \(\operatorname{End}_{\mathbf{A}}(\mathbf{N})\) 的可逆元。假设对所有 \(s \in \mathbf{S}\)，\(h_s\) 在 \(\operatorname{End}_{\mathbf{A}}(\mathbf{N})\) 中可逆；于是元素 \(h_a\)（其中 \(a \in \mathbf{A}\)）与元素 \(h_s\)（其中 \(s \in \mathbf{S}\)）的逆在 \(\operatorname{End}_{\mathbf{A}}(\mathbf{N})\) 中生成一个交换子环 B，且 A 到 B 的同态 \(a \mapsto h_a\) 使 S 的元素的像都可逆。于是由此得出（第 1 节，命题 1），存在唯一的同态 \(h'\)（\(\mathrm{A}[\mathrm{S}^{-1}]\) 到 B），使得

\[
h ^ {\prime} (a / s) = h _ {a} \left(h _ {s}\right) ^ {- 1};
\]

我们知道（《代数学》第 II 章，§ 1，第 14 节），这样的同态在 N 上定义了一个 A{[}S \(^{-1}\) {]}-模结构，使得 (a/s).y = h \(_{s}^{-1}\) (a.y)；由这个 A{[}S \(^{-1}\) {]}-模结构借助同态 i \(_{A}^{S}\) 导出的 A-模结构正是最初给定的结构。

反之，若 N 是 A{[}S \(^{-1}\) {]}-模，并且借助 \(i_{\text{A}}^{\text{S}}\) 把它看作 A-模，则位似映射 \(y \mapsto sy\)（其中 \(s \in \text{S}\)）是双射，因为 \(y \mapsto (1/s)y\) 是 \(y \mapsto sy\) 的逆映射；而由上述程序从 N 的 A-模结构导出的 A{[}S \(^{-1}\) {]}-模结构就是最初给定的 A{[}S \(^{-1}\) {]}-模结构。于是，在 A{[}S \(^{-1}\) {]}-模与由 S 的元素诱导的位似映射均为双射的 A-模之间存在一个典范的一一对应；此外，若 N、N' 是两个具有这一性质的 A-模，则每个 A-模同态 \(u: \text{N} \to \text{N}'\) 也是 N 与 N' 的 A{[}S \(^{-1}\) {]}-模结构之间的同态，因为对所有 \(y \in \text{N}\) 与所有 \(s \in \text{S}\)，可以写出 \(u(y) = u(s. ((1/s)y)) = s.u((1/s)y)\)，由此得 \(u((1/s)y) = (1/s)u(y)\)；反之是显然的。

在这一讨论之下，命题 3 的陈述不过是把纯量扩张到 A{[}S \(^{-1}\) {]} 时由 M 得到的模的特征刻画，其中用到上述解释（《代数学》第 II 章，§ 5，第 1 节，注 1）。

定义 3. 设 A 是环，S 是 A 的子集，\(\overline{\mathbf{S}}\) 是由 S 生成的 A 的乘性子集，M 是 A-模。由 S 定义的 M 的分式模，记作 \(\mathbf{M}[\mathbf{S}^{-1}]\) 或 \(\overline{\mathbf{S}}^{-1}\mathbf{M}\)，是指 \(\mathbf{A}[\mathbf{S}^{-1}]\) -模 \(\mathbf{M} \otimes_{\mathbf{A}} \mathbf{A}[\mathbf{S}^{-1}]\)。

本章中我们通常用 \(i_{\mathbf{M}}^{\mathbf{S}}\) 表示典范映射 \(m \mapsto m \otimes 1\)（\(\mathbf{M}\) 到 \(\mathbf{M}[\mathbf{S}^{-1}]\)）。

注

\begin{enumerate}
\def\labelenumi{(\arabic{enumi})}
\item
  显然 \(\mathbf{M}[\overline{\mathbf{S}}^{-1}] = \mathbf{M}[\mathbf{S}^{-1}]\)。
\item
  对 \(m \in \mathbf{M}\) 与 \(s \in \overline{\mathbf{S}}\)，我们还用 \(m / s\) 表示元素 \(m \otimes (1 / s)\)，它属于 \(\mathbf{M}[\mathbf{S}^{-1}]\)。\(\mathbf{M}[\mathbf{S}^{-1}]\) 的每个元素都具有这种形式，因为这样的元素具有形式 \(\sum_{i} m_i \otimes (a_i / s)\)，其中 \(m_i \in \mathbf{M}\)，\(a_i \in \mathbf{A}\)，\(s \in \mathbf{S}\)（第 1 节，注 2），并且
\[
m _ {i} \otimes (a _ {i} / s) = (a _ {i} m _ {i}) \otimes (1 / s),
\]

于是 \(\sum_{i} m_{i} \otimes (a_{i}/s) = m \otimes (1/s)\)，其中 \(m = \sum_{i} a_{i} m_{i}\)。从而


\[
(m / s) + (m ^ {\prime} / s ^ {\prime}) = (s ^ {\prime} m + s m ^ {\prime}) / s s ^ {\prime}\tag{4}
\]

\[
(a / s) (m / s ^ {\prime}) = (a m) / (s s ^ {\prime})\tag{5}
\]

其中 \(m, m' \in \mathbf{M}, a \in \mathbf{A}\)，\(s, s' \in \mathbf{S}\)。

\item
  若 S 是 A 的素理想 p 的余集，我们就用 \(M_{p}\) 代替 \(S^{-1}M\)。
\item
  设 M 是 A{[}S \(^{-1}\) {]}-模；若把 M 典范地看作 A-模，则 \(i_{M}^{S}\) 是双射，因为由 M 与恒等映射 \(l_{M}\) 组成的有序对显然也是 M{[}S \(^{-1}\) {]} 与 \(i_{M}^{S}\) 所解的泛映射问题的解。于是把 M 与 M{[}S \(^{-1}\) {]} 等同起来。
\end{enumerate}

命题 4. 设 S 是 A 的乘性子集，M 是 A-模。为使 m/s = 0（m ∈ M，s ∈ S），必须且只需存在 s' ∈ S 使得 s'm = 0。

若 \(s' \in S\) 满足 \(s'm = 0\)，则显然 \(m/s = (s'm)/(ss') = 0\)。反之，设 \(m/s = 0\)。由于 \(1/s\) 在 \(S^{-1}A\) 中可逆，有 \(m/1 = 0\)。对每个包含 1 的 A-子模 \(P\)（\(S^{-1}A\) 的子模），我们用 \(\beta(P, m)\) 表示 \((m, 1)\) 在 \(M \times P\) 到 \(M \otimes_A P\) 的典范映射下的像；于是 \(\beta(S^{-1}A, m) = 0\)。我们知道（《代数学》第 II 章，§ 6，第 3 节，命题 7 的推论 4），存在一个包含 1 的有限生成子模 \(P\)（\(S^{-1}A\) 的子模），使得 \(\beta(P, m) = 0\)。对所有 \(t \in S\)，我们用 \(A_t\) 表示形如 \(a/t\)（其中 \(a \in A\)）的元素组成的集合；由于 \(P\) 是有限生成的，存在 \(t \in S\) 使得 \(P \subset A_t\)（第 1 节，注 2），于是 \(\beta(A_t, m) = 0\)。映射 \(a \mapsto a/t\)（从 \(A\) 到 \(A_t\)）是满射；设 \(B\) 是它的核。它定义了一个满射 \(h: M \otimes_A A \to M \otimes_A A_t\)，其核为 BM（把 M 与 \(M \otimes A\) 等同）；于是 \(\beta(A_t, m) = h(tm)\)，从而 \(tm\) 可以表示为形式 \(\sum_{i=1}^{r} b_i m_i\)，其中 \(b_i \in B\)，\(m_i \in M (1 \leqslant i \leqslant r)\) 。由 \(b_i/t = 0\) 可表示为形式 \(\sum_{i=1}^{n} b_i m_i\)，其中 \(b_i \in \mathbf{B}\)，\(m_i \in \mathbf{M}\)（\(1 \leqslant i \leqslant r\)）。由于 \(b_i / t = 0\) 对 \(1 \leqslant i \leqslant r\) 成立，存在 \(t' \in \mathbf{S}\) 使得 \(t'b_i = 0\) 对 \(1 \leqslant i \leqslant r\) 成立，从而 \(t'tm = 0\)，这就证明了命题 4。

推论 1. 为使 \(m / s = m' / s'\)（在 \(S^{-1}M\) 中）成立，必须且只需存在 \(t \in S\) 使得 \(t(s'm - sm') = 0\)。

\[
(m / s) - (m ^ {\prime} / s ^ {\prime}) = (s ^ {\prime} m - s m ^ {\prime}) / s s ^ {\prime}.
\]

推论 2. 设 M 是有限生成的 A-模。为使 \(S^{-1}M = 0\)，必须且只需存在 \(s \in S\) 使得 \(sM = 0\)。

在不对 M 作任何假设的情况下，显然，关系 \(sM = 0\)（对某个 \(s \in S\)）蕴含 \(S^{-1}M = 0\)。反之，设 \(S^{-1}M = 0\)，并设 \((m_i)_{1 \leqslant i \leqslant n}\) 是 M 的一个生成系；\(m_i / 1\) 生成 \(S^{-1}A\)-模 \(S^{-1}M\)，因此说 \(S^{-1}M = 0\) 就相当于说 \(m_i / 1 = 0\) 对 \(1 \leqslant i \leqslant n\) 成立；由

命题 4，存在 \(s_i \in S\) 使得 \(s_i m_i = 0\)；取 \(s = s_1 s_2 \ldots s_n\)，则 \(sm_i = 0\) 对所有 \(i\) 成立，从而 \(sM = 0\)。

推论 3. 设 M 是有限生成的 A-模。为使 A 的理想 a 满足 aM = M，必须且只需存在 \(a \in a\) 使得 \((1 + a)M = 0\)。

显然，关系 \((1 + a)M = 0\) 蕴含 M = aM。为证明其逆，我们使用下面的引理：

引理 1. 对 A 的每个理想 a，由形如 \(1 + a\)（其中 \(a \in a\)）的元素组成的集合 S 是 A 的一个乘性子集，而集合 \(a'\)（由 \(S^{-1}A\) 中形如 a/s 的元素组成，其中 \(a \in a\) 且 \(s \in S\)）是包含于 \(S^{-1}A\) 的 Jacobson 根中的一个理想。

第一个断言是明显的，\(a'\) 是 \(S^{-1}A\) 的理想这一事实也同样明显。另一方面，\((1/1) + (a/s) = (s + a)/s\)，且 \(s + a \in S\) 对所有 \(s \in S\) 与 \(a \in a\)（按 S 的定义）成立；因此 \((1/1) + (a/s)\) 在 \(S^{-1}A\) 中对所有 \(a/s \in a'\) 可逆，这就完成了引理的证明（《代数学》第 VIII 章，§ 6，第 3 节，定理 1）。

在此情形下，若令 \(N = S^{-1}M\)，则显然 N 是有限生成的 \(S^{-1}A\)-模；若 aM = M，则 \(a'N = N\)，于是由 Nakayama 引理（《代数学》第 VIII 章，§ 6，第 3 节，命题 6 的推论）得 N = 0；由此并利用推论 2 即得本推论。

命题 5. 设 A、B 是两个环，S 是 A 的乘性子集，T 是 B 的乘性子集，f 是 A 到 B 的同态且 \(f(S) \subset T\)。设 M 是 A-模，N 是 B-模，u 是 M 到 N 的 A-线性映射。则存在唯一的 \(S^{-1}A\)-线性映射 \(u'\)（从 \(S^{-1}M\) 到 \(T^{-1}N\)），使得 \(u'(m/1) = u(m)/1\) 对所有 \(m \in M\) 成立。

映射 \(i_{\mathbf{N}}^{\mathbf{T}} \circ u\)（从 \(\mathbf{M}\) 到 \(\mathbf{T}^{-1}\mathbf{N}\)）是 A-线性的。此外，若 \(s \in \mathbf{S}\)，则 \(f(s) \in \mathbf{T}\)，因而由 \(s\) 在 \(\mathbf{T}^{-1}\mathbf{N}\) 上诱导的位似映射是双射。于是由命题 3 得出 \(u'\) 的存在性与唯一性。这时，对 \(m \in \mathbf{M}\) 与 \(s \in \mathbf{S}\)，

\[
u ^ {\prime} (m / s) = u (m) / f (s).\tag{6}
\]

沿用同样的记号，设 C 是第三个环，U 是 C 的乘性子集，g 是 B 到 C 的同态且 \(g(T) \subset U\)，P 是 C-模，v 是 N 到 P 的 B-线性映射，\(v'\) 是与 v 相伴的 \(T^{-1}B\)-线性映射（从 B 到 \(U^{-1}P\)）。则

\[
(v \circ u) ^ {\prime} = v ^ {\prime} \circ u ^ {\prime}\tag{7}
\]

其中左边是 A-线性映射 \(S^{-1}M \rightarrow U^{-1}P\)（与 \(v \circ u\) 相伴的）。类似地，若 \(u_{1}\) 是 M 到 N 的第二个 A-线性映射，则

\[
(u + u _ {1}) ^ {\prime} = u ^ {\prime} + u _ {1} ^ {\prime},\tag{8}
\]

左边的项是 A-线性映射 \(S^{-1}M \rightarrow T^{-1}N\)（与 \(u + u_{1}\) 相伴的）。

注 (5). 若在命题 5 中取 B = A、T = S 以及 \(f = 1_{A}\)，则容易看出，\(u'\) 正是映射 \(u \otimes 1: M \otimes S^{-1}A \to N \otimes S^{-1}A\)。今后我们把它记作 \(S^{-1}u\)；若 S 是 A 的素理想 p 的余集，我们就用 \(u_{p}\) 代替 \(S^{-1}u\)。

命题 6. 设 \(f\) 是环 \(\mathbf{A}\) 到环 \(\mathbf{B}\) 的同态，\(\mathbf{S}\) 是 \(\mathbf{A}\) 的乘性子集。存在唯一的映射 \(j\)（从 \((f(\mathbf{S}))^{-1}\mathbf{B}\) 到 \(\mathbf{S}^{-1}\mathbf{B}\)）（其中 \(\mathbf{B}\) 借助 \(f\) 被看作 A-模），使得 \(j(b/f(s)) = b/s\) 对所有 \(b \in \mathbf{B}\)、\(s \in \mathbf{S}\) 成立。若 \(f': \mathbf{S}^{-1}\mathbf{A} \to (f(\mathbf{S}))^{-1}\mathbf{B}\) 是与 \(f\) 相伴的环同态（第 1 节，命题 2），则 \(j \circ f' = \mathbf{S}^{-1}f\)。映射 \(j\) 是 \(\mathbf{S}^{-1}\mathbf{A}\)-模结构（\((f(\mathbf{S}))^{-1}\mathbf{B}\) 上由 \(f'\) 定义者）到 \(\mathbf{S}^{-1}\mathbf{B}\) 上相应结构的同构，也是 B-模结构（\((f(\mathbf{S}))^{-1}\mathbf{B}\) 上者）到 \(\mathbf{S}^{-1}\mathbf{B}\) 上相应结构的同构（这一结构由定义 \(\mathbf{S}^{-1}\mathbf{B} = (\mathbf{S}^{-1}\mathbf{A}) \otimes_{\mathbf{A}} \mathbf{B}\) 而得）。

若 \(b, b'\) 属于 B，\(s, s'\) 属于 S，则条件 \(b / s = b' / s'\) 与 \(b / f(s) = b' / f(s')\) 等价，这由命题 4 的推论 1 得出；这就确立了 \(j\) 的存在性并表明 \(j\) 是双射；\(j\) 的唯一性是明显的。显然 \(j\) 是加群同构。若 \(a \in A\)、\(b \in B\)、\(s \in S\)、\(t \in S\)，则

\[
(a / s). (b / f (t)) = f ^ {\prime} (a / s) (b / f (t)) = f (a) / f (s) (b / f (t)) = (f (a) b) / f (s t),
\]

由此得出 \(j\) 是 \((\mathbf{S}^{-1}\mathbf{A})\)-线性的。显然 \(j \circ f^1 = \mathbf{S}^{-1}f\)。最后，若 \(b \in \mathbf{B}\)、\(b' \in \mathbf{B}\)、\(s \in \mathbf{S}\)，则 \(j(b \cdot (b'/f(s)) = j(bb'/f(s)) = bb'/s = b \cdot (b'/s))\)，这就证明了最后一个断言。

命题 6 中的映射 \(j\) 称为 \((f(\mathbf{S}))^{-1}\mathbf{B}\) 到 \(\mathbf{S}^{-1}\mathbf{B}\) 上的典范同构。这两个集合通常借助 \(f\) 等同起来；这时 \(f' = \mathbf{S}^{-1}f\)，\(i_{\mathbf{B}}^{\mathbf{S}} = i_{\mathbf{B}}^{f(\mathbf{S})}\)。

